% Generado por src/analisis/potencial.py; no editar a mano.
\begin{table}[H]
\centering\small
\caption{Potencia instalable y energía anual estimadas sobre el área conservadora (orden de magnitud). Aprovechamiento del techo 40\,\% (rango 30--50\,\%), 0,20\,kWp/m\textsuperscript{2} de módulo, 4,5\,kWh/m\textsuperscript{2}/día y rendimiento global 0,80.}
\label{tab:potencial}
\begin{tabularx}{\textwidth}{@{}rLrrrrr@{}}
\toprule
 & & \textbf{Área} & \multicolumn{3}{c}{\textbf{Potencia (MWp)}} & \textbf{Energía} \\
\cmidrule(lr){4-6}
\textbf{\#} & \textbf{Municipio} & \textbf{(ha)} & \textbf{Baja} & \textbf{Central} & \textbf{Alta} & \textbf{(GWh/año)} \\
\midrule
1 & Santa Marta & 1\,421 & 853 & 1\,137 & 1\,421 & 1\,494 \\
2 & Valledupar & 1\,351 & 811 & 1\,081 & 1\,351 & 1\,420 \\
3 & Buenaventura & 715 & 429 & 572 & 715 & 751 \\
4 & Florencia & 669 & 401 & 535 & 669 & 703 \\
5 & Turbo & 466 & 280 & 373 & 466 & 490 \\
6 & Santander de Quilichao & 384 & 231 & 307 & 384 & 404 \\
7 & San Andrés de Tumaco & 379 & 228 & 303 & 379 & 399 \\
8 & Ciénaga & 347 & 208 & 278 & 347 & 365 \\
9 & Apartadó & 337 & 202 & 269 & 337 & 354 \\
10 & San José del Guaviare & 310 & 186 & 248 & 310 & 326 \\
11 & Tierralta & 297 & 178 & 238 & 297 & 313 \\
12 & Arauquita & 292 & 175 & 234 & 292 & 307 \\
13 & Tame & 276 & 166 & 221 & 276 & 290 \\
14 & Puerto Asís & 266 & 160 & 213 & 266 & 280 \\
15 & Tibú & 264 & 158 & 211 & 264 & 277 \\
\midrule
\multicolumn{2}{@{}l}{Top 15} & 7\,775 & 4\,665 & 6\,220 & 7\,775 & 8\,173 \\
\multicolumn{2}{@{}l}{170 municipios PDET} & 20\,156 & 12\,093 & 16\,124 & 20\,156 & 21\,188 \\
\bottomrule
\end{tabularx}
\end{table}
